\section{Chapitre~\ref{sec:glutcomputer}. Ce que calculent les neurones \nt{GLUT}}

Les énoncés logiques du chapitre sont des théorèmes sur les circuits à poids positifs ou nuls, démontrés dans le chapitre lui-même; l'expérience confirme le modèle de neurone et le fait que les circuits construits d'après ces théorèmes (détecteur de coïncidences, ligne à retard, groupe auto-entretenu, chaîne) se rencontrent réellement dans le cerveau.

{\footnotesize
\setlength{\tabcolsep}{3pt}\renewcommand{\arraystretch}{1.15}
\begin{longtable}{>{\raggedright}p{0.5cm}>{\raggedright}p{4.0cm}>{\raggedright}p{5.6cm}>{\raggedright}p{2.3cm}>{\raggedright\arraybackslash}p{1.7cm}}
\toprule
N\textsuperscript{o} & Affirmation & Expérience et ce qui est mesuré & Source & Statut \\
\midrule\endfirsthead
\toprule
N\textsuperscript{o} & Affirmation & Expérience et ce qui est mesuré & Source & Statut \\
\midrule\endhead
1 & Le neurone est un circuit RC avec seuil et remise à zéro~\eqref{eq:glutneuron} & On a injecté dans le corps de neurones pyramidaux du cortex du rat (tranches) un courant bruité semblable au flux synaptique du cerveau vivant. La dépendance de la fréquence moyenne à la moyenne et à la dispersion du courant est décrite par un neurone intégrateur à fuite, si l'on ajoute une lente adaptation de la fréquence. Les plages de $\tau$ et des retards sont données dans le chapitre sans référence expérimentale & \cite{Rauch2003} & mesuré \\
2 & Le modèle~\eqref{eq:glutneuron} est une simplification: les branches d'entrée émettent elles-mêmes des impulsions locales & Enregistrement directement sur les dendrites de neurones pyramidaux du cortex visuel de la souris in vivo: un stimulus visuel déclenche des impulsions dendritiques locales dépendant des récepteurs NMDA; les supprimer élargit l'accord du neurone sur l'orientation & \cite{Smith2013} & mesuré \\
3 & Fenêtre de coïncidence $\Delta_{\max}=\tau\ln\frac{w}{\theta-w}$~\eqref{eq:window}; pour $\theta=1.5w$, elle vaut $0.7\tau$ & Conséquence du modèle~\eqref{eq:glutneuron}. Une mesure directe d'une fenêtre de sommation étroite existe pour les neurones à inhibition rapide (chapitre~\ref{sec:gaba}, ligne~3 de son tableau) & déduction du chapitre & théorie \\
4 & Un réseau fait seulement de neurones \nt{GLUT} ne calcule que des fonctions monotones des entrées (proposition~\ref{prop:monotone}) & Démonstration dans le chapitre: itération à partir du silence avec un membre de droite non décroissant. C'est un énoncé sur le modèle; aucune expérience ne l'a vérifié sur un réseau vivant privé d'inhibition & déduction du chapitre & théorie \\
5 & Les organes des sens fournissent une paire de fils: certaines cellules répondent au renforcement du stimulus, d'autres à son affaiblissement & Rétine: dès les cellules bipolaires, le signal des cônes se sépare en canaux ON et OFF. Blocage pharmacologique des cellules bipolaires ON chez le singe: les canaux restent séparés jusqu'au cortex; après le blocage, la détection d'une augmentation de lumière se dégrade, mais pas celle d'une diminution & \cite{Schiller1992} & mesuré \\
6 & Avec le code à deux fils, un réseau de neurones \nt{GLUT} calcule n'importe quelle fonction logique de façon asynchrone, sans cadence commune, au prix de deux fois plus de neurones (proposition~\ref{prop:dualrail}, \eqref{eq:dualrail}, fig.~\ref{fig:dualxor}) & Le code à deux fils et la détection de fin de calcul par le signal lui-même sont une technique standard des circuits asynchrones insensibles aux retards. Le circuit OU exclusif à six neurones est construit dans le chapitre. Aucune expérience ne montre que le cerveau calcule les fonctions logiques précisément en code à deux fils & \cite{Sparso2001} & théorie \\
7 & La plus grande différence de temps d'arrivée du son aux deux oreilles est chez l'homme de $\approx0.6\ms$, et le cerveau distingue environ dix microsecondes & Auditeurs dans une expérience à choix forcé entre deux réponses: seuil de discrimination de la différence de temps (75\,\% de réponses correctes) de $9$\,\textmu s pour un bruit dans la bande $150$--$1700$\,Hz, $11$\,\textmu s pour un son pur de $1$\,kHz, $28$\,\textmu s pour un clic. La limite de $0.6\ms$ est un calcul du chapitre, $0.22\,\text{m}/343\,\text{m/s}$ & \cite{KlumppEady1956} & mesuré \\
8 & Chez les oiseaux, la différence de temps se change en position grâce à des lignes à retard et des détecteurs de coïncidences~\eqref{eq:itd} (fig.~\ref{fig:itd}) & Effraie des clochers: les axones venus des deux oreilles entrent dans le noyau des détecteurs de coïncidences par des côtés opposés; le temps d'arrivée des impulsions le long de l'axone varie de façon ordonnée avec la profondeur, et l'étendue des retards à travers le noyau ($160$\,\textmu s sur $700$\,\textmu m) correspond à la plage des retards interauraux & \cite{Carr1990} & mesuré \\
9 & Chez les mammifères, on n'a pas trouvé de rangée de retards; la même tâche est résolue par des détecteurs de coïncidences avec une inhibition précisément calée venue de neurones \nt{GLYC} & Enregistrement in vivo dans l'olive supérieure médiane de la gerbille: les réponses ne forment pas de carte des directions; l'application de glycine et de son bloqueur par micropipette montre que l'inhibition \emph{glycinergique} précisément calée dans le temps fixe la plage utile des retards. L'inhibition vient ici de \nt{GLYC}, non de \nt{GABA} & \cite{Brand2002,Grothe2010} & mesuré \\
10 & Un groupe à forte rétroaction est une bascule; l'activité auto-entretenue persiste des secondes après le stimulus (fig.~\ref{fig:bistable}) & Cortex préfrontal du singe dans une tâche de saccade différée vers un point mémorisé (délai de $1$--$6$\,s): $87$ des $170$ neurones liés à la tâche changent de fréquence pendant le délai, et pour $79\,\%$ d'entre eux cela dépend de la direction du point mémorisé. Que cette activité soit entretenue par une rétroaction à deux états stables relève de la théorie & \cite{Funahashi1989,Wang2001} & mesuré \\
11 & Le bit s'écrit si l'apport dure plus de $\approx0.7\tau$ (fig.~\ref{fig:recurrentgroup}b) & Calcul de l'équation $\tau\,\dd r/\dd t=-r+f(wr+I)$ par la méthode d'Euler pour $w=1$, $I=0.6$; pas de script dans \texttt{models/}. Pas d'expérience & calcul du chapitre & modèle \\
12 & La mémoire peut être gardée dans la facilitation des synapses, presque sans impulsions & Théorie: le calcium résiduel dans les terminaisons garde l'information environ une seconde. Appui: dans le cortex préfrontal médian du furet (enregistrement de plusieurs neurones à la fois), il existe un sous-réseau de neurones pyramidaux reliés par des synapses facilitatrices à long effet rémanent. Revue des observations d'intervalles \og silencieux\fg{} pendant le maintien en mémoire. Quelle façon le cortex utilise et quand reste une question ouverte & \cite{Mongillo2008,WangMarkram2006,Stokes2015} & hypothèse \\
13 & La mémoire s'efface par fatigue: la réserve de neurotransmetteur s'épuise & Enregistrements appariés de neurones pyramidaux du cortex: la réponse de la synapse baisse lors d'impulsions fréquentes, la vitesse de la baisse est fixée par la probabilité de libération. Qu'elle éteigne précisément un groupe auto-entretenu n'est pas montré directement & \cite{Tsodyks1997} & mesuré \\
14 & Une chaîne de groupes est une minuterie déclenchée par un événement; chez les oiseaux chanteurs, les neurones se déclenchent strictement à tour de rôle & Enregistrement de neurones identifiés du centre vocal supérieur du diamant mandarin pendant le sommeil et le chant: chaque neurone qui projette vers le noyau moteur émet une seule brève bouffée à un instant précis du chant, et les neurones se déclenchent successivement & \cite{Hahnloser2002} & mesuré \\
\bottomrule
\end{longtable}}
